% Generated draft: requires mathematical and source review.
\documentclass[11pt]{article}
\usepackage[margin=25mm]{geometry}
\usepackage{fontspec}
\setmainfont{Latin Modern Roman}
\usepackage{amsmath,amssymb,mathtools,unicode-math}
\setmathfont{Latin Modern Math}
\usepackage{enumitem,needspace,xurl,hyperref}
\setlength{\parindent}{0pt}
\setlength{\parskip}{5pt}
\setlength{\emergencystretch}{3em}
\newcommand{\N}{\mathbb{N}}
\newcommand{\Z}{\mathbb{Z}}
\newcommand{\Q}{\mathbb{Q}}
\newcommand{\R}{\mathbb{R}}
\newcommand{\C}{\mathbb{C}}
\newcommand{\F}{\mathbb{F}}
\newcommand{\A}{\mathbb{A}}
\newcommand{\PP}{\mathbb{P}}
\newcommand{\E}{\mathbb{E}}
\newcommand{\eps}{\varepsilon}
\newcommand{\dd}{\,\mathrm{d}}
\begin{document}
\section*{UnsolvedMath ID: 11000224}
\textbf{Source problem code:} catalogue entry
\subsection{Definitions and mathematical statement}
For a translation surface $X$, consider its directional straight-line flows away from the conical singularities. Write $\mathcal M(X)$ for minimal directions, $\mathcal U(X)$ for uniquely ergodic directions, and $\mathcal P(X)$ for completely periodic directions. Conventionally, minimality means density of the nonsingular flow orbits in the surface, unique ergodicity means uniqueness of the normalized invariant transverse measure for that directional flow, and complete periodicity means decomposition into parallel cylinders outside the critical graph. Singular fixed-point measures are not included in this standard directional-flow interpretation; precise saddle-connection conventions must be checked.

Problem 17 asks to characterize surfaces satisfying each of
\[
 (1)\quad\mathcal M(X)=\mathcal U(X),\qquad
 (2)\quad\mathcal P(X)=\Theta\setminus\mathcal U(X),
\]
where $\Theta$ is the set of directions (oriented or modulo reversal, consistently). The source states that (2) implies (1).

The source reports that in genus $2$ every non-Veech surface has a minimal but non-uniquely-ergodic direction [ChMa]. It also reports an example due to B. Weiss, using Hubert--Schmidt [HuSt2], that is not Veech but satisfies (2): a ramified cover of a Veech surface with a single ramification point. The example's genus and the distinction between a ramification point and a branch value are not supplied. It is not asserted to be a genus-two example. These claims rule out simply replacing the requested characterization by ``Veech surfaces'' without the appropriate genus and property distinctions.

The following extracted counting discussion is context. A nonsingular closed orbit determines a parallel cylinder, whose circumference is the orbit length. The source gives for the standard flat torus the asymptotic $\pi L^2/\zeta(2)$ for cylinders of length at most $L$, and reports quadratic asymptotics for Veech surfaces [Ve3] and generic surfaces in each stratum [EsMa]. Counting conventions matter: counting primitive lattice vectors counts both orientations, while identifying opposite vectors introduces a factor of two. Maximal-cylinder/family conventions, area normalization and the meaning of generic must be recovered before adopting that coefficient. No new cylinder-counting question is added to the two characterization requests.
\subsection{Short English statement}
Characterize translation surfaces where every minimal direction is uniquely ergodic, and the stronger class where non-uniquely-ergodic directions are exactly completely periodic ones. Preserve the genus-two result and the reported higher-scope non-Veech example.
\subsection{Sources}
\begin{itemize}
\item Local registry: \url{data/unsolvedmath/problems.json}, permanent UnsolvedMath ID $11000224$, exact source code \texttt{catalogue entry}.
\item Farb, \emph{Problems on Mapping Class Groups and Related Topics}, Problem 17, PDF page 262, extracted chapter footer P. Hubert, H. Masur, T. Schmidt and A. Zorich: \url{https://www.math.uchicago.edu/~farb/papers/mcgbook.pdf}. Supplied access date: 2026-07-29; extraction: pdf-font-structure; page-disambiguated numbering among 329 headings.
\item Supplied citations: [ChMa] for the genus-two assertion; Hubert--Schmidt [HuSt2] and an example credited to B. Weiss; [Ve3] for Veech-surface counting; [EsMa] for generic-stratum counting. Full bibliography entries are absent. The registry cannot verify the modern status.
\end{itemize}
\end{document}
